% part: intuitionistic-logic
% chapter: soundness-completeness
% section: canonical-model

\documentclass[../../../include/open-logic-section]{subfiles}

\begin{document}

\olfileid{int}{sc}{mod}

\olsection{કૅનોનિકલ નિદર્શ}

આપણા નિદર્શનાં જગતો પ્રાકૃતિક સંખ્યાઓની સાન્ત
શ્રેણીઓ~$\sigma$ હશે, એટલે $\sigma \in \Nat^*$.
નોંધો કે $\Nat^*$ને આગમનાત્મક રીતે આ પ્રમાણે વ્યાખ્યાયિત
કરીએ છીએ:
\begin{enumerate}
\item $\emptyseq \in \Nat^*$.
\item જો $\sigma \in \Nat^*$ અને $n \in \Nat$ હોય,
  તો $\sigma.n \in \Nat^*$ (અહીં $\sigma.n$ એ
  $\sigma \concat \tuple{n}$ છે અને $\sigma \concat \sigma'$
  એ $\sigma$ તથા $\sigma'$નું જોડાણ છે).
\item $\Nat^*$માં બીજું કશું નથી.
\end{enumerate}
તેથી આગમનાત્મક વ્યાખ્યાઓ આપવા $\Nat^*$ વાપરી શકીએ છીએ.

$\tuple{!B_1, !C_1}$, $\tuple{!B_2, !C_2}$, \dots,
એ !!{formula}sની બધી જોડીઓની પરિગણના હોય.
!!{formula}sનો ગણ~$\Delta$ આપેલો હોય, તો
$\Delta(\sigma)$ને નીચે મુજબ આગમનાત્મક રીતે
વ્યાખ્યાયિત કરીએ છીએ:
\begin{enumerate}
\item $\Delta(\emptyseq) = \Delta$
\item $\Delta(\sigma.n) = {}$
  \[
  \begin{cases}
    (\Delta(\sigma) \cup \{!B_n\})^* &
    \text{જો $\Delta(\sigma) \cup \{!B_n\} \Proves/ !C_n$ હોય} \\
    \Delta(\sigma) & \text{અન્યથા}
  \end{cases}
  \]
\end{enumerate}
અહીં $(\Delta(\sigma) \cup \{!B_n\})^*$નો અર્થ
!!{formula}sનો એવો અભાજ્ય ગણ છે, જે
$\Delta(\sigma) \cup \{!B_n\}$ અને !!{formula}~$!C_n$
પર \olref[lin]{lem:lindenbaum} લાગુ કરતાં મળે છે.
આ વ્યાખ્યા મુજબ જો $\Delta(\sigma) \cup \{!B_n\}
\Proves/ !C_n$ હોય, તો $\Delta(\sigma.n)
\Proves !B_n$ અને $\Delta(\sigma.n) \Proves/ !C_n$.
એ પણ નોંધો કે કોઈ પણ~$n$ માટે
$\Delta(\sigma) \subseteq \Delta(\sigma.n)$.
જો $\Delta$ અભાજ્ય હોય, તો દરેક~$\sigma$ માટે
$\Delta(\sigma)$ અભાજ્ય હોય છે.

\begin{defn}\ollabel{defn:canonical-model}
  ધારો કે $\Delta$ અભાજ્ય છે. ત્યારે $\Delta$ માટેનું
  \emph{કૅનોનિકલ નિદર્શ} $\mModel{M(\Delta)}$ આ પ્રમાણે
  વ્યાખ્યાયિત થાય છે:
  \begin{enumerate}
  \item $W = \Nat^*$, એટલે પ્રાકૃતિક સંખ્યાઓની સાન્ત
    શ્રેણીઓનો ગણ.
  \item $R$ એવો આંશિક ક્રમ છે કે $R\sigma\sigma'$
    ત્યારે અને માત્ર ત્યારે જ જ્યારે $\sigma$ એ~$\sigma'$નો
    પ્રારંભિક ખંડ હોય (એટલે કોઈ શ્રેણી~$\sigma''$ માટે
    $\sigma' = \sigma \concat \sigma''$ હોય).
  \item $V(p) = \Setabs{\sigma}{p \in \Delta(\sigma)}$.
  \end{enumerate}
\end{defn}

$R$ ખરેખર આંશિક ક્રમ છે તે સહેલાઈથી ચકાસી શકાય છે.
$V$ માટેની એકદિશવર્ધિતાની શરત પણ સંતોષાય છે.
$\Delta(\sigma) \subseteq \Delta(\sigma.n)$ હોવાથી
$R\sigma\sigma'$ હોય ત્યારે $\sigma$ પછી જોડાતા ભાગની
લંબાઈ પર આગમન કરીને
$\Delta(\sigma) \subseteq \Delta(\sigma')$ મળે છે.
\sourcecorrection{OLMD-095}{મૂળ આ છેલ્લું આગમન
$\sigma$ પર કહે છે; $\sigma$ સ્થિર રાખીને તેના પછી
જોડાતી સાન્ત શ્રેણીની લંબાઈ પર આગમન જોઈએ.}

\end{document}
